% !TeX root = ../../portafoliofinal.tex
% !TeX encoding = utf8
% ====================================================================
%  CAPÍTULO — INFORME GLOBAL DE LA EMPRESA
%
%  Análisis integrado de la gestión de RRHH en CaixaBank y
%  recomendaciones transversales que interrelacionan las distintas
%  prácticas estudiadas a lo largo del curso.
% ====================================================================

\chapter{Informe global de la gestión de Recursos Humanos en CaixaBank}

\section{Introducción: la gestión de RRHH como sistema integrado}

A lo largo de las tres prácticas previas, el equipo ha analizado de forma
detallada las principales prácticas de recursos humanos de CaixaBank:
el análisis y diseño de puestos, la planificación de plantilla, el
reclutamiento, la selección, la integración, la formación, el desarrollo
de carrera, la evaluación del rendimiento y la retribución.

El propósito de este informe global no es repetir ni resumir lo ya expuesto,
sino adoptar una mirada sistémica que permita identificar las \textbf{sinergias
y dependencias entre las distintas prácticas} y formular recomendaciones que,
al interrelacionar varios ámbitos de RRHH simultáneamente, generen un impacto
transformador mayor que el que cada mejora aislada podría producir.

CaixaBank es una organización extraordinariamente compleja: más de 45.000
empleados, una red de más de 4.000 oficinas en España, presencia internacional
y una transformación digital en marcha que afecta transversalmente a todos
los procesos de gestión del talento. En este contexto, las políticas de RRHH
no pueden diseñarse en silos; deben formar un ecosistema coherente en el que
cada decisión sobre un área refuerce o complemente las demás.

% ================================================================
\section{Análisis integrado: fortalezas y zonas de mejora detectadas}
% ================================================================

\subsection{Fortalezas sistémicas del modelo de RRHH de CaixaBank}

Del análisis conjunto de las tres prácticas se desprende que CaixaBank
presenta un modelo de gestión de recursos humanos que destaca en varios
ámbitos interrelacionados:

\begin{itemize}
  \item \textbf{Ecosistema tecnológico integrado:} La plataforma SAP
        SuccessFactors actúa como columna vertebral que conecta el análisis
        de puestos, la planificación, el reclutamiento, la evaluación del
        desempeño y la retribución. La integración de Microsoft Copilot añade
        una capa de inteligencia artificial que potencia la analítica predictiva
        en todas estas áreas.

  \item \textbf{Coherencia entre planificación y reclutamiento:} La entidad
        gestiona el déficit de perfiles tecnológicos —identificado en el Tema~3—
        mediante una estrategia de reclutamiento mixto que combina la movilidad
        interna (detallada en el Tema~4) con programas específicos de captación
        de talento joven (\textit{CaixaBank Talent Program}).

  \item \textbf{Modelo retributivo competitivo y transparente:} La estructura
        de retribución —analizada en el Tema~8— supera entre un 30 y un 40\,\%
        los mínimos del IV Convenio Colectivo de Banca, incluye un plan de
        retribución flexible y publica sus indicadores de equidad en el informe
        de sostenibilidad, lo que refuerza tanto la atracción como la retención
        del talento.

  \item \textbf{Cultura de aprendizaje continuo:} La plataforma Evoluciona
        con más de 45.000 usuarios activos, el modelo 70-20-10, los programas
        de coaching y mentoring y las alianzas con instituciones académicas
        configuran un ecosistema formativo maduro y ampliamente implantado.
\end{itemize}

\subsection{Zonas de mejora y tensiones sistémicas identificadas}

Sin embargo, el análisis transversal revela también ciertas tensiones o
incoherencias entre prácticas que merecen atención:

\begin{itemize}
  \item \textbf{Tensión entre digitalización y humanización:} CaixaBank
        avanza decididamente hacia procesos de RRHH impulsados por IA
        (reclutamiento, evaluación, formación), pero esta automatización
        genera el riesgo de despersonalizar la relación empleado-empresa.
        Este riesgo fue explícitamente señalado por la directora de sucursal
        entrevistada al referirse al proceso de selección.

  \item \textbf{Presión comercial y evaluación del rendimiento:} Los sindicatos
        han denunciado que los sistemas de \textit{ranking} y control diario
        generan tensión y riesgo de \textit{burnout}. Esta presión contrasta
        con la vocación declarada de CaixaBank como «banca responsable» y
        con sus métricas ASG, creando una incoherencia entre el discurso
        institucional y la experiencia real del empleado de red.

  \item \textbf{Desconexión entre \textit{outplacement} y gestión de la
        salida:} Tal como se constató en la práctica de Planificación (Tema~3),
        CaixaBank no dispone de un programa permanente de recolocación; el
        \textit{outplacement} se aplicó exclusivamente durante el ERE de 2021.
        Esta carencia contrasta con la solidez del proceso de atracción e
        integración, generando una asimetría en el ciclo de vida del empleado.

  \item \textbf{Evaluación del rendimiento principalmente descendente:} A pesar
        de las tendencias hacia la evaluación 360° señaladas en el Tema~7,
        la práctica real sigue siendo mayoritariamente descendente. Esto limita
        la detección de problemas de liderazgo y cultura en los niveles medios
        de la organización.
\end{itemize}

% ================================================================
\section{Recomendaciones globales e integradoras}
% ================================================================

Las siguientes recomendaciones no se limitan a un área aislada, sino que
proponen intervenciones que conectan múltiples prácticas de RRHH con el
objetivo de generar un impacto sistémico y sostenible.

\begin{recomendacion}{1}{Ecosistema de datos de talento «Talent Intelligence»}
\textbf{Práticas que interrelaciona:} Planificación (T3) + Reclutamiento (T4) + Evaluación del
rendimiento (T7) + Formación (T6).

\medskip
\textbf{Descripción:}
CaixaBank dispone de un volumen de datos sobre sus empleados sin precedentes
en el sector, distribuidos entre SAP SuccessFactors, Evoluciona y los módulos
de People Analytics. La recomendación consiste en crear un panel de
\textbf{\textit{Talent Intelligence}} que integre en tiempo real:
\begin{itemize}
  \item Los datos de desempeño individual (SAP SuccessFactors).
  \item El historial formativo y progreso de cada empleado en Evoluciona.
  \item Las brechas competenciales detectadas en la evaluación periódica.
  \item Las vacantes abiertas en el portal interno de movilidad.
\end{itemize}
Este panel permitiría, por ejemplo, que cuando se abre una vacante interna
(Tema~4), el sistema proponga automáticamente los tres empleados con mayor
ajuste competencial, considerando no solo su puesto actual sino su evolución
formativa reciente. Asimismo, si la evaluación del rendimiento (Tema~7)
detecta brechas en un empleado, el sistema generaría automáticamente un
itinerario formativo personalizado en Evoluciona (Tema~6) sin necesidad de
intervención manual del \textit{Business Partner} de RRHH.

\medskip
\textbf{Impacto esperado:}
Reducción del \textit{Time to Fill} interno, mayor personalización del
desarrollo, y una planificación de plantilla más precisa al anticipar qué
empleados están listos para asumir nuevas responsabilidades.
\end{recomendacion}

\begin{recomendacion}{2}{\textit{Journey} del empleado integrado: de la socialización al desarrollo}
\textbf{Prácticas que interrelaciona:} Selección e integración (T5) + Formación (T6) + Desarrollo de carrera (T7).

\medskip
\textbf{Descripción:}
Actualmente, el proceso de \textit{onboarding} (Tema~5) y el plan de formación
inicial (Tema~6) existen como programas paralelos que no están suficientemente
integrados con el plan de desarrollo de carrera a largo plazo (Tema~7).
La propuesta es diseñar un \textbf{\textit{Employee Journey} integrado} desde
el primer día de trabajo hasta el desarrollo profesional pleno, con tres fases
conectadas:
\begin{enumerate}
  \item \textbf{Fase de acogida} (meses 1--3): Socialización formal e informal,
        asignación de \textit{buddy}, acceso guiado a Evoluciona y primera
        conversación de desarrollo con el \textit{Business Partner}.
  \item \textbf{Fase de consolidación} (meses 4--12): Completar los itinerarios
        normativos y técnicos; primera evaluación de desempeño; definición del
        plan de carrera inicial basado en fortalezas y aspiraciones.
  \item \textbf{Fase de desarrollo} (año 2 en adelante): Rotación de puestos,
        acceso a programas de talento, coaching ejecutivo y revisión anual del
        plan de carrera alineada con la evaluación de desempeño.
\end{enumerate}
Este diseño secuencial no es solo un cambio de procesos; exige que los
responsables directos estén formados y comprometidos con su rol como
\textit{gestores del talento}, y que la evaluación del desempeño (Tema~7) deje
de ser un trámite anual para convertirse en un hito estructural del
\textit{Employee Journey}.

\medskip
\textbf{Nueva tendencia aplicada:}
El concepto de \textbf{\textit{Talent Experience Management}} (TXM), que fusiona
la tecnología LMS con el CRM del empleado, está siendo adoptado por entidades
como ING y BBVA para sustituir los procesos fragmentados por una experiencia
continua y personalizada del ciclo de vida laboral.
\end{recomendacion}

\begin{recomendacion}{3}{Cierre del ciclo de vida --- \textit{Outplacement} estructural y gestión de la salida}
\textbf{Prácticas que interrelaciona:} Planificación sustractiva (T3) + Formación (T6) + Retribución
(T8) + Marca empleadora (T4).

\medskip
\textbf{Descripción:}
Una de las incoherencias más llamativas del modelo de RRHH de CaixaBank
es la asimetría entre la entrada y la salida. La entidad dedica recursos
considerables a atraer, integrar y desarrollar a sus empleados, pero no
dispone de un programa permanente de \textit{outplacement} que asista
a quienes salen en circunstancias distintas a una prejubilación negociada.

La propuesta consiste en establecer un \textbf{Programa de Transición
Profesional Permanente} que incluya:
\begin{itemize}
  \item \textbf{Módulos de empleabilidad en Evoluciona:} Accesibles para
        cualquier empleado que inicie un proceso de salida voluntaria o
        involuntaria, con contenidos sobre marca personal, búsqueda de empleo
        y actualización de competencias.
  \item \textbf{Alianzas con \textit{outplacement} externo:} Para casos de
        reducción de plantilla individual no vinculados a un ERE colectivo,
        activar convenios con empresas especializadas de recolocación sin
        coste adicional para el empleado.
  \item \textbf{Entrevistas de salida sistematizadas:} El Director de RRHH
        entrevistado confirmó que se realizan cuando «el empleado se presta».
        La propuesta es convertirlas en obligatorias y anónimas, canalizado
        su análisis mediante People Analytics para detectar patrones de abandono
        y corregirlos proactivamente.
\end{itemize}
\textbf{Vinculación con la retribución (Tema~8):}
Parte del «salario emocional» que CaixaBank ofrece debería incluir explícitamente
la garantía de un proceso de salida digno y asistido, como elemento diferenciador
de la propuesta de valor al empleado (\textit{EVP}) frente a otros empleadores
del sector.

\medskip
\textbf{Impacto esperado:}
Mejora de la percepción de CaixaBank como empleador responsable,
reducción del «síndrome del superviviente» entre quienes permanecen,
y generación de embajadores de marca entre quienes salen con una
experiencia positiva.
\end{recomendacion}

\begin{recomendacion}{4}{Humanización de la evaluación del rendimiento mediante el modelo 360\textdegree{} y la desvinculación del \textit{ranking}}
\textbf{Prácticas que interrelaciona:} Evaluación del rendimiento (T7) + Formación (T6) + Retribución
variable (T8) + Clima laboral.

\medskip
\textbf{Descripción:}
El modelo actual de evaluación del rendimiento en CaixaBank combina datos
objetivos de producción comercial con una evaluación descendente del
responsable directo. Esta configuración, sumada a la existencia de
\textit{rankings} y reportes diarios de resultados, genera tensiones
documentadas por los sindicatos y el riesgo de primar los resultados
a corto plazo sobre la calidad del servicio y el bienestar del empleado.

La recomendación propone una evolución del sistema en dos ejes:
\begin{enumerate}
  \item \textbf{Evaluación 360°:} Incorporar, al menos para perfiles de mando
        intermedio y directivo, la posibilidad de que los miembros del equipo
        evalúen de forma confidencial a su responsable. Esto permitiría al
        departamento de Personas detectar estilos de liderazgo tóxicos que
        generan resultados a costa de dañar el clima del equipo.
  \item \textbf{Desvinculación parcial del \textit{ranking} de los objetivos
        individuales:} Sustituir progresivamente los objetivos puramente
        individuales por una combinación de metas individuales (50\,\%),
        metas de equipo de oficina (30\,\%) y métricas de calidad de
        servicio y bienestar del cliente (20\,\%). Esta combinación reduce
        la competencia destructiva entre compañeros y alinea mejor el
        modelo de incentivos con la visión de «banca sostenible» declarada
        en el Plan Estratégico 2025--2027.
\end{enumerate}
\textbf{Vinculación con la formación (Tema~6) y la retribución (Tema~8):}
Los resultados de la evaluación 360° deben alimentar directamente los
itinerarios de formación en liderazgo disponibles en Evoluciona, cerrando
el ciclo diagnóstico-aprendizaje. Asimismo, la incorporación de métricas
ASG en el bonus —ya iniciada para la alta dirección— debería extenderse
progresivamente a los mandos intermedios.

\medskip
\textbf{Nueva tendencia aplicada:}
El concepto de \textbf{\textit{Performance Enablement}} (habilitación del
rendimiento), popularizado por empresas como Adobe y Deloitte, propone
sustituir las revisiones anuales por conversaciones continuas y estructuradas
orientadas al desarrollo. Varios bancos europeos (ING, Erste Group) ya lo
han implementado con resultados positivos en retención y clima laboral.
\end{recomendacion}

\begin{recomendacion}{5}{\textit{Reskilling} masivo como política de planificación preventiva}
\textbf{Prácticas que interrelaciona:} Planificación de RRHH (T3) + Formación (T6) + Desarrollo de
carrera (T7) + Procesos sustractivos (T3).

\medskip
\textbf{Descripción:}
El Plan Estratégico 2025--2027 de CaixaBank identifica un déficit de perfiles
tecnológicos que se está cubriendo mediante la contratación de 3.000 nuevos
empleados. Sin embargo, la entidad cuenta con miles de empleados de red cuyo
perfil —orientado a la banca tradicional de ventanilla— puede verse gradualmente
desplazado por la digitalización.

La propuesta es articular un \textbf{Programa de \textit{Reskilling} Preventivo y
Continuo} que, a diferencia de los programas formativos reactivos actuales, funcione
como una política estructural de planificación:
\begin{itemize}
  \item \textbf{Identificación temprana mediante People Analytics:} Usar SAP
        SuccessFactors para modelizar qué perfiles tienen mayor probabilidad de
        quedar obsoletos en un horizonte de 3--5 años y proponerles itinerarios
        de reconversión antes de que la brecha sea irreversible.
  \item \textbf{Itinerarios de transición ``perfil analógico → perfil híbrido'':}
        Diseñar rutas formativas en Evoluciona que permitan a un empleado de
        ventanilla con 15 años de experiencia convertirse en gestor de banca digital
        o especialista en atención omnicanal, reconociendo y aprovechando su
        conocimiento del negocio como ventaja diferencial.
  \item \textbf{Colaboración con universidades:} Ampliar los actuales acuerdos
        con la UPF Barcelona School of Management y otras instituciones para incluir
        módulos de certificación en análisis de datos, ciberseguridad básica
        e inteligencia artificial aplicada a las finanzas.
\end{itemize}
\textbf{Impacto en la planificación sustractiva (Tema~3):}
Un programa de \textit{reskilling} bien ejecutado reduce la necesidad de recurrir
a EREs o prejubilaciones anticipadas para ajustar la plantilla, lo que supone
un ahorro económico significativo (el ERE de 2021 costó 1.900 millones de euros)
y un impacto social positivo que refuerza la reputación de CaixaBank como empleador
responsable.
\end{recomendacion}

\begin{recomendacion}{6}{Integración de la estrategia ASG en el ciclo completo de RRHH}
\textbf{Prácticas que interrelaciona:} Análisis de puestos (T2) + Selección (T5) + Evaluación del
rendimiento (T7) + Retribución (T8) + \textit{Employer branding} (T4).

\medskip
\textbf{Descripción:}
CaixaBank ha integrado métricas ASG (Ambientales, Sociales y de Gobernanza) en
la retribución variable de la alta dirección y publica compromisos de diversidad
e inclusión en su informe de sostenibilidad. Sin embargo, esta integración no
ha permeado de manera consistente en los procesos operativos de RRHH.

La recomendación consiste en extender el compromiso ASG a todos los eslabones
del ciclo de vida del empleado:
\begin{itemize}
  \item \textbf{Análisis de puestos con sesgos auditados (T2):} Revisar las
        descripciones de puestos para eliminar sesgos de género o edad en
        las especificaciones de candidatos. CaixaBank ya dispone de una
        «Oficina de Inteligencia Artificial» que podría incorporar esta función.
  \item \textbf{Selección con criterios DEI medibles (T5):} Establecer
        indicadores de diversidad (género, edad, discapacidad) en el cuadro de
        mando del departamento de Personas, vinculados al rendimiento de los
        equipos de \textit{Talent Acquisition}.
  \item \textbf{Evaluación del rendimiento con competencias de inclusión (T7):}
        Incorporar la competencia «gestión inclusiva de equipos» en la evaluación
        de todos los puestos de mando, con un peso específico en la puntuación
        final que genere un incentivo real para los directivos.
  \item \textbf{Retribución variable con componente ASG para mandos intermedios (T8):}
        Extender el modelo ya vigente para la alta dirección a los directores
        de zona y de oficina, vinculando una parte del bonus a indicadores
        de bienestar del equipo (tasa de absentismo, satisfacción de empleado,
        tasa de bajas por estrés).
\end{itemize}
\textbf{Nueva tendencia aplicada:}
El concepto de \textbf{\textit{HR Sustainability}} propone integrar los objetivos
de sostenibilidad en todos los procesos de gestión del capital humano, no solo
en los de nivel directivo. El \textit{Responsible Business Alliance} identifica
esta integración como uno de los principales diferenciadores de las empresas
que logran mantener el compromiso de sus empleados a largo plazo.
\end{recomendacion}

% ================================================================
\section{Limitaciones generales en la obtención de información}
% ================================================================

El equipo considera importante reconocer de forma explícita las limitaciones
que han condicionado el alcance del análisis a lo largo de las tres entregas:

\begin{enumerate}

  \item \textbf{Confidencialidad de datos internos.}
        En las tres entrevistas realizadas, los interlocutores de CaixaBank
        —Director de RRHH regional y Directora de sucursal— no pudieron proporcionar
        cifras presupuestarias específicas sobre reclutamiento, formación o
        retribución variable, amparándose en la política de confidencialidad de
        la entidad. Ello ha obligado al equipo a recurrir a fuentes secundarias
        (Informe Anual, Plan Estratégico, informe de sostenibilidad) o a
        proponer estimaciones razonadas basadas en el sector.

  \item \textbf{Perspectiva limitada de la fuente.}
        Las entrevistas se realizaron con responsables de RRHH a nivel regional
        o de sucursal. Esta perspectiva aporta información valiosa sobre los
        procesos operativos, pero no permite acceder a las decisiones estratégicas
        que se toman en la sede central (Dirección de Personas y Organización).
        Algunas respuestas fueron de alcance limitado («no supo responder»),
        lo que indica que el entrevistado no tenía visibilidad sobre determinados
        aspectos corporativos.

  \item \textbf{Desfase temporal de las fuentes.}
        Parte de la información proviene de documentos públicos (Plan Estratégico
        2025--2027, Informe Anual 2024) que reflejan la situación de la entidad
        en el momento de su publicación, pero que pueden no capturar cambios
        organizativos más recientes. CaixaBank es una empresa en continua
        transformación, y algunos datos (número de empleados, red de oficinas,
        métricas de formación) pueden haber variado.

  \item \textbf{Ausencia de perspectiva del empleado de base.}
        Las fuentes utilizadas son predominantemente institucionales (informes
        corporativos, entrevistas con responsables). La perspectiva de los
        empleados de base —que sería fundamental para evaluar la coherencia
        entre el discurso de CaixaBank y la experiencia real de trabajo— se ha
        podido incorporar solo de forma parcial a través de plataformas como
        Glassdoor, con las limitaciones de representatividad que esto implica.

\end{enumerate}

% ================================================================
\section{Conclusión: hacia un modelo de RRHH integrado y sostenible}
% ================================================================

CaixaBank es, sin duda, una de las empresas españolas con mayor sofisticación
en la gestión de sus recursos humanos. Dispone de tecnología de vanguardia,
un ecosistema formativo maduro y una política retributiva competitiva. Sin
embargo, el análisis integral realizado en este portafolio pone de manifiesto
que la verdadera ventaja competitiva sostenible no reside en cada práctica
aislada, sino en la capacidad de integrarlas en un sistema coherente y
autoreforzado.

Las seis recomendaciones presentadas en este informe persiguen precisamente
ese objetivo: cerrar los círculos abiertos entre prácticas, eliminar las
inconsistencias entre el discurso institucional y la experiencia real del
empleado, y preparar a la organización para los desafíos de una banca cada
vez más digital y más humana al mismo tiempo.

La transformación que el sector bancario experimenta no es solo tecnológica;
es una transformación del contrato psicológico entre empresa y empleado.
CaixaBank tiene los recursos, la escala y la voluntad declarada para liderar
esa transformación también en el ámbito del capital humano.

\endinput
